\documentclass[../DoAn.tex]{subfiles}
\begin{document}

Chương này trình bày các giải pháp và đóng góp nổi bật của đồ án. Các nội dung được lựa chọn là những phần có ảnh hưởng trực tiếp tới khả năng chơi được của sản phẩm: xử lý từ tiếng Việt, sinh bảng ô chữ, tổ chức màn chơi theo chủ đề và lưu tiến độ người chơi.

\section{Xử lý từ tiếng Việt trong trò chơi ô chữ}
Vấn đề đầu tiên của đồ án là biểu diễn từ tiếng Việt sao cho vừa phù hợp với bảng ô chữ, vừa giữ được dạng đáp án tự nhiên cho người chơi. Một đáp án như một cụm từ tiếng Việt có thể chứa khoảng trắng, trong khi bảng ô chữ cần xếp từng ký tự liên tiếp. Ngoài ra, các chữ cái có dấu cần được giữ nguyên để người chơi nhìn thấy đúng đáp án.

Giải pháp của đồ án là tách mỗi từ thành hai dạng. Dạng thứ nhất là từ đầy đủ sau khi chuẩn hóa chữ thường, dùng để hiển thị và so khớp đáp án. Dạng thứ hai là chuỗi xếp bảng, được tạo bằng cách loại bỏ khoảng trắng. Mỗi ký tự trong chuỗi xếp bảng được chuyển thành thông tin gồm ký tự gốc, ký tự cơ sở, biến thể nguyên âm và thanh điệu. Nhờ đó, hệ thống có thể đặt đúng chữ có dấu trên bảng, đồng thời chấp nhận đáp án người chơi nhập theo dạng có hoặc không có khoảng trắng.

Kết quả của giải pháp này là trò chơi có thể xử lý các từ đơn và cụm từ tiếng Việt trong cùng một cơ chế. Người chơi nhập đáp án theo cách tự nhiên, trong khi bộ sinh bảng vẫn làm việc với chuỗi ký tự liên tục.

Điểm đáng chú ý của giải pháp là phần xử lý tiếng Việt không bị trộn trực tiếp vào thuật toán sinh bảng. Bộ sinh bảng chỉ cần biết một \texttt{WordEntry} có danh sách ký tự và độ dài tương ứng. Việc tạo danh sách ký tự, xử lý dấu và chuẩn hóa đáp án được thực hiện trước đó bởi các lớp tiện ích. Điều này giúp thuật toán sinh bảng có thể tập trung vào bài toán không gian hai chiều, trong khi logic ngôn ngữ được gom vào một nhóm thành phần riêng.

Một lợi ích khác là hệ thống có khả năng mở rộng cách so khớp đáp án trong tương lai. Hiện tại, đáp án được chấp nhận khi khớp đầy đủ hoặc khớp sau khi bỏ khoảng trắng. Nếu cần hỗ trợ thêm chế độ nhập không dấu hoặc chế độ luyện tập dễ hơn, có thể bổ sung vào lớp chuẩn hóa mà không cần thay đổi toàn bộ giao diện hoặc bộ sinh bảng. Đây là một ví dụ cho việc tách trách nhiệm giúp mã nguồn dễ phát triển tiếp.

\begin{table}[H]
\centering
\begin{tabularx}{\textwidth}{|>{\RaggedRight\arraybackslash}p{3.3cm}|>{\RaggedRight\arraybackslash}p{4.1cm}|>{\RaggedRight\arraybackslash}X|}
\hline
\textbf{Dạng dữ liệu} & \textbf{Ví dụ} & \textbf{Mục đích sử dụng} \\ \hline
Từ hiển thị & \texttt{bánh trung thu} & Dùng trong đáp án, lịch sử đoán và dữ liệu nội dung. \\ \hline
Dạng xếp bảng & \texttt{bánhtrungthu} & Dùng để tính số ô và đặt ký tự lên bảng. \\ \hline
Danh sách ký tự & \texttt{b, á, n, h, ...} & Dùng để so khớp giao điểm giữa các từ. \\ \hline
Thông tin ký tự & Ký tự gốc, ký tự cơ sở, biến thể, thanh điệu & Dùng để lưu đúng chữ tiếng Việt trong từng ô. \\ \hline
\end{tabularx}
\caption{Các dạng biểu diễn từ tiếng Việt trong hệ thống}
\label{table:vietnamese_word_representation}
\end{table}

\section{Sinh bảng ô chữ tự động}
Nếu tạo thủ công từng bảng ô chữ, việc mở rộng số lượng chủ đề và màn chơi sẽ tốn nhiều thời gian. Vì vậy, đồ án xây dựng bộ sinh bảng tự động từ danh sách từ của mỗi màn. Thách thức của bài toán là phải đặt được nhiều từ vào bảng, bảo đảm các từ không chồng sai ký tự, không nối liền không hợp lệ và có giao điểm để tạo cảm giác ô chữ.

Bộ sinh bảng sử dụng chiến lược đặt từ theo từng bước. Từ đầu tiên được đặt nằm ngang gần trung tâm. Với mỗi từ tiếp theo, hệ thống tìm các ký tự trùng với những từ đã đặt, thử đặt theo hướng vuông góc và kiểm tra điều kiện hợp lệ. Mỗi vị trí hợp lệ được chấm điểm dựa trên số giao điểm, độ gần trung tâm và mức tăng diện tích bao của toàn bộ bảng chữ. Khi không tìm được vị trí theo cách chính, hệ thống dùng bước tìm kiếm dự phòng trên toàn bảng nhưng vẫn yêu cầu có giao điểm.

Giải pháp này không phải thuật toán tối ưu toàn cục, nhưng phù hợp với quy mô ứng dụng vì thời gian sinh nhanh, mã nguồn dễ hiểu và chất lượng bảng đủ dùng cho các màn chơi nhỏ. Đây là lựa chọn cân bằng giữa độ phức tạp thuật toán và yêu cầu thực tế của sản phẩm.

Trong thiết kế ban đầu của nhiều bài toán ô chữ, một phương án có thể là thử quay lui toàn bộ để tìm cách đặt được nhiều từ nhất. Tuy nhiên, với danh sách từ tiếng Việt theo chủ đề, số lượng từ có thể khá lớn và độ dài từ không đồng đều. Nếu áp dụng quay lui đầy đủ, thời gian sinh có thể tăng nhanh và gây cảm giác chờ khi người chơi bắt đầu màn. Đồ án vì vậy chọn hướng heuristic nhiều lần thử. Mỗi lần thử thay đổi thứ tự từ theo một seed ổn định, sau đó lấy kết quả đặt được nhiều từ nhất. Cách này tăng khả năng tìm được bảng tốt hơn một lần thử duy nhất nhưng vẫn giữ quá trình sinh trong giới hạn chấp nhận được.

Một quyết định quan trọng là loại bỏ các từ cô lập sau khi sinh. Trong trò chơi ô chữ, các từ nên có liên kết với nhau để người chơi có thể suy luận từ giao điểm. Nếu một từ được đặt mà không giao với từ nào khác, nó giống một câu hỏi riêng biệt hơn là một phần của bảng ô chữ. Vì vậy, sau khi đặt từ, hệ thống kiểm tra các từ không có giao điểm và dựng lại bảng chỉ với các từ liên kết. Điều này có thể làm giảm số từ cuối cùng của màn, nhưng giúp chất lượng bảng phù hợp hơn với mục tiêu trò chơi.

\begin{table}[H]
\centering
\begin{tabularx}{\textwidth}{|>{\RaggedRight\arraybackslash}p{3.3cm}|>{\RaggedRight\arraybackslash}X|>{\RaggedRight\arraybackslash}X|}
\hline
\textbf{Tiêu chí} & \textbf{Cách đánh giá trong thuật toán} & \textbf{Ý nghĩa} \\ \hline
Số giao điểm & Đếm số ô đã có ký tự trùng với từ mới & Tăng tính liên kết giữa các từ trong bảng \\ \hline
Độ gần trung tâm & Tính khoảng cách điểm giữa của từ tới tâm bảng & Giúp bảng không lệch quá xa một phía \\ \hline
Diện tích bao & Tính phần diện tích tăng thêm sau khi đặt từ & Ưu tiên bảng gọn và dễ hiển thị \\ \hline
Tính hợp lệ & Kiểm tra biên bảng, ký tự xung đột và ô kề & Tránh tạo từ sai hoặc va chạm không hợp lệ \\ \hline
\end{tabularx}
\caption{Các tiêu chí đánh giá vị trí đặt từ}
\label{table:placement_scoring}
\end{table}

\subsection{Đánh giá bộ sinh bảng}
\label{subsection:5.2.1}
Để đánh giá bộ sinh bảng, đồ án chạy một kịch bản đo tự động ở chế độ không giao diện (\code{--headless}) trên Godot 4.6.3. Kịch bản lần lượt sinh bảng cho toàn bộ màn của 28 chủ đề, ghi lại số từ hợp lệ đưa vào mỗi màn, số từ thực sự đặt được sau khi loại từ cô lập và thời gian sinh. Mỗi chủ đề được chia thành 3 màn, tạo thành 84 màn được đo.

Kết quả tổng hợp cho thấy cả 84/84 màn đều sinh được bảng hợp lệ đủ điều kiện chơi (đạt ngưỡng tối thiểu \code{MIN\_STAGE\_WORDS}). Tổng cộng có 767 từ hợp lệ được đưa vào, trong đó 649 từ được đặt lên bảng, tương ứng tỉ lệ đặt từ trung bình 84,6\%. Xét theo từng màn, tỉ lệ đặt từ dao động từ 57,1\% đến 100\%; xét theo từng chủ đề, tỉ lệ nằm trong khoảng 68,4\% đến 100\%. Thời gian sinh tổng cộng là 261,57\,ms cho 84 màn, trung bình khoảng 3,1\,ms mỗi màn, đủ nhanh để không gây cảm giác chờ khi người chơi bắt đầu màn.

Phần từ không được đặt chủ yếu là các từ không tìm được giao điểm phù hợp với phần bảng đang dựng và bị loại trong bước bỏ từ cô lập. Đây là sự đánh đổi có chủ đích: hệ thống ưu tiên một bảng gắn kết, các từ giao nhau để người chơi suy luận được, thay vì cố đặt tối đa số từ và chấp nhận các từ rời rạc.

\begin{longtable}{|>{\RaggedRight\arraybackslash}p{3.6cm}|>{\Centering\arraybackslash}p{2.2cm}|>{\Centering\arraybackslash}p{2.2cm}|>{\Centering\arraybackslash}p{2.2cm}|>{\Centering\arraybackslash}p{2.2cm}|}
\caption{Kết quả đo bộ sinh bảng trên 28 chủ đề (Godot 4.6.3, chế độ headless)}
\label{table:board_benchmark}\\
\hline
\textbf{Chủ đề} & \textbf{Từ đưa vào} & \textbf{Từ đặt được} & \textbf{Tỉ lệ (\%)} & \textbf{Thời gian (ms)} \\ \hline
\endfirsthead
\hline
\textbf{Chủ đề} & \textbf{Từ đưa vào} & \textbf{Từ đặt được} & \textbf{Tỉ lệ (\%)} & \textbf{Thời gian (ms)} \\ \hline
\endhead
Phòng ngủ & 25 & 24 & 96,0 & 4,22 \\ \hline
Thời gian & 22 & 19 & 86,4 & 6,70 \\ \hline
Nhà tắm & 25 & 22 & 88,0 & 6,64 \\ \hline
Các kiểu tóc & 20 & 18 & 90,0 & 4,68 \\ \hline
Trò chơi dân gian & 28 & 21 & 75,0 & 13,50 \\ \hline
Vũ khí & 32 & 28 & 87,5 & 9,80 \\ \hline
Trung thu & 20 & 16 & 80,0 & 8,42 \\ \hline
Phòng khách & 24 & 24 & 100,0 & 2,22 \\ \hline
Giáng sinh & 18 & 18 & 100,0 & 1,81 \\ \hline
Thiên tai & 17 & 15 & 88,2 & 5,53 \\ \hline
Đám cưới & 26 & 23 & 88,5 & 7,25 \\ \hline
Đại dương & 48 & 39 & 81,3 & 16,56 \\ \hline
Dụng cụ y tế & 20 & 19 & 95,0 & 3,51 \\ \hline
Nhạc cụ & 30 & 25 & 83,3 & 10,13 \\ \hline
Cơ quan nội tạng & 20 & 18 & 90,0 & 5,31 \\ \hline
Năm mới & 32 & 26 & 81,3 & 11,10 \\ \hline
Nông trại & 41 & 39 & 95,1 & 7,07 \\ \hline
Hình dáng & 16 & 14 & 87,5 & 5,56 \\ \hline
Khuôn mặt & 21 & 19 & 90,5 & 4,64 \\ \hline
Tính cách & 47 & 35 & 74,5 & 22,94 \\ \hline
Đồ uống & 25 & 25 & 100,0 & 2,46 \\ \hline
Tình yêu & 21 & 18 & 85,7 & 7,01 \\ \hline
Thể thao & 45 & 38 & 84,4 & 13,59 \\ \hline
Côn trùng & 19 & 13 & 68,4 & 20,93 \\ \hline
Vũ trụ & 31 & 24 & 77,4 & 13,49 \\ \hline
Môn học & 18 & 13 & 72,2 & 11,47 \\ \hline
Âm nhạc & 37 & 28 & 75,7 & 17,70 \\ \hline
Phim ảnh & 39 & 28 & 71,8 & 17,33 \\ \hline
\textbf{Tổng cộng} & \textbf{767} & \textbf{649} & \textbf{84,6} & \textbf{261,57} \\ \hline
\end{longtable}

\section{Tổ chức nội dung theo chủ đề và màn chơi}
Một yêu cầu quan trọng của trò chơi là người chơi cần có cảm giác tiến triển thay vì chỉ giải một bảng đơn lẻ. Đồ án giải quyết yêu cầu này bằng cách tổ chức dữ liệu thành các chủ đề như phòng ngủ, thời gian, nhà tắm, trò chơi dân gian, nhạc cụ và nhiều nhóm từ khác. Mỗi chủ đề được chia thành một số màn, mỗi màn chứa một tập từ đủ nhỏ để bảng không quá dày và người chơi có thể hoàn thành trong thời gian hợp lý.

Việc chia màn được thực hiện tự động từ danh sách từ của chủ đề. Các từ được sắp xếp ưu tiên theo độ dài ở lần thử đầu tiên, sau đó phân bổ vào các nhóm màn. Khi người chơi chọn một màn, hệ thống mới sinh bảng cho màn đó. Cách làm này giúp giảm thời gian xử lý khi vào menu, đồng thời làm cho quá trình chơi theo chủ đề rõ ràng hơn.

Cách tổ chức theo chủ đề cũng giúp nội dung gợi ý có ngữ cảnh hơn. Cùng một từ có thể mang nghĩa khác nhau trong các chủ đề khác nhau, vì vậy hệ thống hỗ trợ gợi ý đặc thù theo chủ đề bên cạnh gợi ý mặc định. Ví dụ, một từ xuất hiện trong chủ đề lễ hội có thể cần gợi ý khác với khi xuất hiện trong chủ đề không gian hoặc đồ vật. Việc gắn chủ đề cho từ giúp câu gợi ý tự nhiên hơn và giảm khả năng gây nhập nhằng.

Một vấn đề khác là độ dài của mỗi màn. Nếu một màn có quá ít từ, trò chơi kết thúc quá nhanh và bảng thiếu liên kết. Nếu một màn có quá nhiều từ, bảng có thể khó sinh, giao diện nhiều gợi ý và người chơi dễ mệt. Đồ án đặt ngưỡng tối thiểu cho số từ trong một màn và giới hạn kích thước bảng, từ đó buộc dữ liệu nội dung phải được chia thành các cụm vừa phải. Cách làm này chưa phải cơ chế cân bằng độ khó hoàn chỉnh, nhưng là nền tảng để phát triển hệ thống độ khó trong tương lai.

Ở phiên bản hiện tại, dữ liệu nội dung gồm 28 chủ đề với tổng cộng 780 từ, trung bình khoảng 27,9 từ mỗi chủ đề. Chủ đề nhỏ nhất là ``Hình dáng'' với 16 từ, chủ đề lớn nhất là ``Đại dương'' với 50 từ. Sau bước chuẩn hóa, loại trùng và lọc theo độ dài hợp lệ, có 767 từ được đưa vào quá trình sinh bảng. Bảng~\ref{table:theme_word_counts} liệt kê một số chủ đề tiêu biểu cùng số từ tương ứng.

\begin{table}[H]
\centering
\small
\begin{tabularx}{\textwidth}{|>{\RaggedRight\arraybackslash}p{4.5cm}|>{\Centering\arraybackslash}p{2.5cm}|>{\RaggedRight\arraybackslash}X|}
\hline
\textbf{Chủ đề} & \textbf{Số từ} & \textbf{Ngữ cảnh nội dung} \\ \hline
Đại dương & 50 & Biển cả, sinh vật nước mặn và hoạt động ven bờ \\ \hline
Tính cách & 47 & Đặc điểm cư xử, thái độ và khí chất con người \\ \hline
Thể thao & 45 & Thi đấu, vận động cơ thể và kỹ năng rèn luyện \\ \hline
Nông trại & 44 & Chăn nuôi, trồng trọt và lao động ở vùng quê \\ \hline
Phim ảnh & 39 & Điện ảnh, thể loại phim và quá trình sản xuất \\ \hline
Vũ khí & 34 & Công cụ tấn công, phòng thủ hoặc gây sát thương \\ \hline
Năm mới & 32 & Ngày tết, phong tục, may mắn và sum họp gia đình \\ \hline
Phòng ngủ & 25 & Nghỉ ngơi, đồ đạc cá nhân và không gian riêng tư \\ \hline
Trung thu & 20 & Đêm rằm, trăng, lễ hội và ký ức tuổi thơ \\ \hline
Hình dáng & 16 & Hình học, đường nét và dạng khối cơ bản \\ \hline
\multicolumn{2}{|l|}{\textbf{Tổng: 28 chủ đề / 780 từ}} & \textit{Trung bình 27,9 từ/chủ đề} \\ \hline
\end{tabularx}
\caption{Một số chủ đề tiêu biểu và số từ tương ứng}
\label{table:theme_word_counts}
\end{table}

\section{Lưu tiến độ và cơ chế sao}
Để tăng tính liên tục của trò chơi, đồ án bổ sung hệ thống tài khoản cục bộ và lưu tiến độ theo từng tài khoản. Mỗi tài khoản có tổng số sao, điểm cao nhất, số lượt chơi và tiến độ từng màn. Khi hoàn thành màn, hệ thống ghi nhận số sao cao nhất, điểm cao nhất và thời gian tốt nhất.

Cơ chế sao được dùng cho cả phần thưởng và chi phí gợi ý. Người chơi có một lượng sao ban đầu, có thể dùng sao để mở một chữ hoặc mở cả từ. Khi hoàn thành màn, người chơi nhận sao dựa trên các nhiệm vụ như hoàn thành màn, đạt điểm mục tiêu, hạn chế số gợi ý, đạt chuỗi đúng hoặc hoàn thành trong thời gian yêu cầu. Cơ chế này tạo thêm mục tiêu phụ ngoài việc đoán hết từ, giúp màn chơi có động lực lặp lại.

Hệ thống nhiệm vụ được sinh theo từng màn dựa trên trạng thái của màn chơi. Một nhiệm vụ luôn là hoàn thành màn, một nhiệm vụ liên quan tới điểm mục tiêu, và nhiệm vụ còn lại được chọn trong một nhóm nhiệm vụ phụ. Nhóm nhiệm vụ phụ có thể là giới hạn số lần dùng gợi ý, đoán đúng liên tiếp, hoàn thành trong thời gian nhất định, không dùng gợi ý mở cả từ hoặc đạt tỉ lệ đúng tối thiểu. Cách thiết kế này làm cho mỗi lượt chơi có thêm tiêu chí đánh giá ngoài kết quả đúng/sai.

Việc dùng sao cho gợi ý cũng tạo ra một vòng lặp tài nguyên đơn giản. Người chơi có thể dùng sao để vượt qua màn khó, nhưng nếu dùng quá nhiều gợi ý thì khó đạt một số nhiệm vụ phụ. Khi hoàn thành màn tốt, người chơi nhận thêm sao để dùng cho các màn sau. Vòng lặp này không phức tạp như hệ thống kinh tế trong các trò chơi lớn, nhưng đủ để tạo quyết định nhỏ cho người chơi: tự giải để giữ sao hay dùng gợi ý để hoàn thành màn.

\begin{table}[H]
\centering
\begin{tabularx}{\textwidth}{|>{\RaggedRight\arraybackslash}p{3.8cm}|>{\RaggedRight\arraybackslash}p{4.2cm}|>{\RaggedRight\arraybackslash}X|}
\hline
\textbf{Loại nhiệm vụ} & \textbf{Điều kiện} & \textbf{Tác dụng trong trải nghiệm} \\ \hline
Hoàn thành màn & Tất cả từ trong màn được giải đúng & Bảo đảm mục tiêu cơ bản của màn chơi \\ \hline
Điểm mục tiêu & Điểm đạt mức yêu cầu & Khuyến khích giải nhanh và hạn chế phụ thuộc gợi ý \\ \hline
Giới hạn gợi ý & Số gợi ý đã dùng không vượt ngưỡng & Khuyến khích người chơi tự suy luận \\ \hline
Chuỗi đúng & Đạt số từ đúng liên tiếp theo yêu cầu & Khuyến khích nhập đáp án cẩn thận \\ \hline
Thời gian & Hoàn thành trong thời gian mục tiêu & Tạo thử thách tốc độ cho màn đã quen \\ \hline
Độ chính xác & Tỉ lệ đáp án đúng đạt ngưỡng & Hạn chế đoán thử quá nhiều \\ \hline
\end{tabularx}
\caption{Các nhóm nhiệm vụ trong màn chơi}
\label{table:mission_types}
\end{table}

\section{Kết chương}
Các giải pháp trên tạo nên đóng góp chính của đồ án: một quy trình xử lý từ tiếng Việt phù hợp với trò chơi ô chữ, một bộ sinh bảng tự động đủ đơn giản để bảo trì, một mô hình tổ chức chủ đề--màn chơi rõ ràng và một cơ chế lưu tiến độ cục bộ. Những đóng góp này giúp sản phẩm không chỉ là một giao diện nhập đáp án, mà là một trò chơi hoàn chỉnh ở mức nguyên mẫu.

Các đóng góp này cũng được kiểm chứng bằng số liệu đo thực tế. Với 28 chủ đề và 780 từ, bộ sinh bảng tạo được 84/84 màn hợp lệ (tỉ lệ 100\%), đặt được 649/767 từ hợp lệ (84,6\%) với thời gian sinh trung bình khoảng 3,1\,ms mỗi màn (xem Bảng~\ref{table:board_benchmark}). Bên cạnh đó, 14 trường hợp kiểm thử chức năng đều đạt trong đợt chạy kiểm thử tự động (xem Mục~\ref{section:4.7}). Những số liệu này cho thấy các giải pháp đề xuất hoạt động ổn định trong phạm vi nguyên mẫu.

\end{document}
